\documentclass[14pt,a4paper]{extarticle}
\usepackage[T2A]{fontenc}
\usepackage[utf8]{inputenc}
\usepackage[english,russian]{babel}
\usepackage{lmodern}
\usepackage{geometry}
\usepackage{array}
\usepackage{tabularx}
\usepackage{longtable}
\usepackage{ltablex}
\usepackage{booktabs}
\usepackage{hyperref}
\usepackage{indentfirst}
\usepackage{microtype}
\usepackage{graphicx}
\usepackage{setspace}
\usepackage{listings}
\usepackage{graphicx}
\usepackage{caption}
\usepackage{float}
\usepackage{tocloft}
\usepackage{titlesec}
\usepackage{newtxtext,newtxmath}

\titleformat{\section}        % форматируем \section
  {\normalfont\bfseries}      % стиль: обычный + жирный
  {\thesection.}              % номер секции с точкой
  {2mm}                       % расстояние между номером и текстом
  {}                           % перед заголовком можно добавить что-то
\titlespacing*{\section}{1.5cm}{*3}{*1}  % {отступ слева}{отступ сверху}{отступ снизу}

\captionsetup{justification=centering}
\renewcommand{\ttdefault}{cmtt}
\keepXColumns
\setlength{\parindent}{3em}
\newgeometry{left=2.5cm,right=2.5cm,top=3cm,bottom=2.5cm}

\renewcommand{\contentsname}{\hspace*{4em}{\smallСодержание}}
\setlength{\cftbeforesecskip}{0.5em}
\renewcommand{\cftsecleader}{\cftdotfill{\cftdotsep}}
\renewcommand{\cftsecaftersnum}{.\hspace{2mm}}

\title{Отчёт по культорологии}
\author{Выполнил студент группы М-РИТ-25-1, Станиславчук Сергей}

\begin{document}

\setstretch{1.5}
\noindent
\begin{minipage}[t]{0.01\textwidth}
    \includegraphics{~/programmer/++Programmer/5_1/LSTU.jpg}
\end{minipage}%
\begin{minipage}[c]{1.1\linewidth}
    \centering
    {\small
    МИНИСТЕРСТВО НАУКИ И ВЫСШЕГО ОБРАЗОВАНИЯ РФ \\[1em]
    ФЕДЕРАЛЬНОЕ ГОСУДАРСТВЕННОЕ БЮДЖЕТНОЕ \\
    ОБРАЗОВАТЕЛЬНОЕ УЧРЕЖДЕНИЕ ВЫСШЕГО ОБРАЗОВАНИЯ \\
    «ЛИПЕЦКИЙ ГОСУДАРСТВЕННЫЙ ТЕХНИЧЕСКИЙ УНИВЕРСИТЕТ»
    }
\end{minipage}%

\vspace{3em}

\noindent
\hspace{1em}%
\begin{minipage}[t]{0.8\textwidth}
\textbf{Институт} \hspace{2em} компьютерных наук \\
\textbf{Кафедра} \hspace{2em} автоматизированных систем управления
\end{minipage}


\vspace{6em}

\begin{center}
\textbf{ЗАДАНИЕ 7} \\
по дисциплине \\ <<Межкультурное взаимодействие>>
\end{center}

\vspace{6em}

\noindent
\begin{tabular*}{\textwidth}{@{\extracolsep{\fill}} l c r }
    Студент М-РИТ-25-1 & \makebox[4cm]{\hrulefill} & Станиславчук С.М. \\[2em]
Руководитель & \\
Ст. преподаватель & \makebox[4cm]{\hrulefill} & Болдырихин О.В. \\
\end{tabular*}

\vfill
\begin{center}
    Липецк, 2025
\end{center}

\thispagestyle{empty}

\newpage
\tableofcontents
\newpage

\section{Задание 1. Вербальные и невербальные коммуникации и их особенности в различных странах}
\begin{sloppypar}
Вербальная коммуникация опирается на речь, словарный запас, интонацию и
    структуру высказываний. Она выражает смысл буквально, но даже здесь
    существуют культурные отличия. В странах с прямым стилем общения —
    например, Германии, США или Нидерландах — ценится точность формулировок,
    чёткие высказывания и открытое выражение позиции. В культурах с непрямой
    коммуникацией, таких как Япония, Китай или арабские государства,
    значительная часть смысла передаётся намёками, контекстом и подтекстом, а
    прямота может восприниматься как грубость.

    Невербальная коммуникация включает жесты, мимику, позы, дистанцию,
    тактильность, зрительный контакт и интонацию голоса. Она варьируется ещё
    сильнее. Одни и те же жесты могут иметь противоположные значения в разных
    культурах: жест «ок» в одних странах нейтрален, а в других считается
    оскорбительным. В Южной Европе допустимо активное использование жестов и
    ближняя дистанция общения, тогда как в странах Северной Европы и Японии
    предпочитают сдержанные движения и большую личную дистанцию.

    Важны также правила зрительного контакта. В США и Европе прямой взгляд
    сигнализирует открытость и уверенность, а в некоторых восточных культурах —
    неуважение или излишнюю напористость. Тактильность тоже различается:
    латиноамериканские и южноевропейские культуры нормально воспринимают частые
    прикосновения, тогда как скандинавские или британские общества относятся к
    этому сдержанно.
\end{sloppypar}

\section{Задание 2. Условия эффективной межкультурной коммуникации.}
\begin{sloppypar}
Эффективная межкультурная коммуникация основывается на понимании различий между
    культурами и умении адаптировать своё поведение к новому контексту. Первое
    условие — осведомлённость о культурных нормах и ценностях собеседника. Это
    включает знание особенностей этикета, предпочтительных форм общения,
    допустимой дистанции, роли жестов и особенностей речевых стратегий. Без
    этого велик риск неверно интерпретировать слова и поведение другого
    человека.

    Второе условие — развитие навыков эмпатии и умения смотреть на ситуацию
    глазами представителя другой культуры. Это позволяет избегать поспешных
    выводов, стереотипов и этноцентризма. Способность услышать, уточнить и
    прояснить смысл помогает сглаживать возможные барьеры.

    Третье условие — гибкость и готовность адаптироваться. Человек, умеющий
    корректировать тон общения, стиль аргументации и невербальное поведение в
    зависимости от культурного контекста, быстрее достигает взаимопонимания.
    Ригидность, наоборот, приводит к конфликтам и недопониманию.

    Четвёртое условие — владение базовыми навыками межкультурной рефлексии.
    Важно осознавать собственные культурные привычки и особенности и понимать,
    что они не универсальны. Это помогает видеть потенциальные источники
    напряжения и вовремя корректировать коммуникацию.
\end{sloppypar}

\section{Задание 3. Влияние культуры на процесс переговоров}
\begin{sloppypar}
Культура напрямую определяет стиль ведения переговоров, способы аргументации,
    правила поведения и саму логику взаимодействия. Люди из разных культурных
    сред подходят к переговорам по-разному, потому что их нормы общения, модели
    принятия решений и отношение к конфликтам существенно различаются.

    Первое, что влияет на переговоры, — предпочтение прямого или непрямого
    стиля общения. Представители стран с прямой коммуникацией (Германия, США,
    Скандинавия) выражают позицию открыто, требуют ясности и быстро переходят к
    обсуждению сути. В культурах непрямого общения (Япония, Корея, арабские
    страны) значительная часть смысла передаётся через контекст, жесты, паузы,
    а прямые высказывания могут восприниматься как давление.

    Второй фактор — отношение ко времени. В одних культурах ценят точность,
    планирование и быстрые решения. В других — ориентируются на длительное
    выстраивание доверия, где сроки гибкие, а скорость уступает важности
    межличностных отношений. Непонимание этих различий легко приводит к
    напряжению на переговорах.

    Третий аспект — иерархия и роль статуса. В странах с высоким уровнем
    дистанции власти решение обычно принимает старший по статусу участник,
    поэтому без него процесс фактически не двигается. В странах с низкой
    дистанцией власти участники чувствуют себя свободнее, решения коллективные,
    а статус не играет столь резкой роли.

    Не меньшее значение имеют нормы невербального поведения — дистанция,
    взгляд, допустимые жесты, особенности этикета. Неосознанное нарушение
    культурных правил может создать недоверие, даже если вербально всё звучит
    корректно.
\end{sloppypar}

\section{Задание 4. Модели ведения переговоров}
\begin{sloppypar}
1. Жёсткая (позиционная) модель переговоров Эта модель основана на чётком
    отстаивании заранее сформулированной позиции. Каждая сторона стремится
    получить максимальную выгоду, не демонстрируя гибкости. Главные инструменты
    — давление, настойчивость, контроль над ситуацией. Переговоры превращаются
    в соревнование, где уступки рассматриваются как слабость. Такая модель
    часто приводит к тупиковым ситуациям или односторонним решениям, которые
    плохо переживают дальнейшую реализацию.

    2. Мягкая позиционная модель В этой модели стороны ориентированы на
    поддержание отношений и избегание конфликта. Уступки делаются легко, ради
    сохранения гармонии. Главная цель — не победа, а предотвращение обострения.
    Такая модель удобна в краткосрочной перспективе, но часто ведёт к
    неравномерному распределению выгод и накоплению скрытого недовольства.

    3. Принципиальная (интересно-ориентированная) модель переговоров Подход,
    предложенный исследователями Гарвардского проекта переговоров. Он
    фокусируется не на позициях, а на реальных интересах сторон. Вместо борьбы
    за формулировки стороны ищут взаимовыгодное решение, анализируют причины
    требований и пытаются найти общий базис. Здесь важны объективные критерии,
    открытый обмен информацией и совместный поиск вариантов. Это одна из самых
    конструктивных и эффективных моделей.

    4. Интегративная модель переговоров Цель — создать ценность, а не делить
    её. Стороны стремятся расширить общую «зону выигрыша», найти новые ресурсы
    или способы сотрудничества. Такая модель строится на доверии, аналитике и
    готовности рассматривать нестандартные варианты. Используется в
    бизнес-партнёрствах, где обе стороны выиграют больше от сотрудничества, чем
    от конфликта.

    5. Дистрибутивная модель переговоров Основана на принципе распределения
    фиксированного ресурса: чем больше получает одна сторона, тем меньше
    остаётся другой. Применяется в ситуациях купли-продажи, финансовых сделок
    или трудовых конфликтов. Здесь важна тактика торга, умение скрывать
    информацию и правильно выбирать момент уступок.

    6. Кооперативная модель Стороны изначально рассматривают друг друга как
    партнёров, а не соперников. Они открыто делятся информацией, совместно
    формируют цели и договариваются о единых правилах взаимодействия. Эта
    модель эффективна там, где отношения важнее результата или когда
    сотрудничество длительное.

    7. Конкурентная модель Близка к жёсткой, но ещё более напористая. Главная
    цель — выиграть любой ценой. Используются тактики давления, якорения,
    «взятия на измор». Подходит для краткосрочных сделок, но высок риск
    разрушить отношения и вызвать ответную агрессию.

    8. Модель с акцентом на построение отношений Типична для восточных и
    традиционных культур. Переговоры начинаются с установления доверия, общения
    о личном, ритуалов вежливости. Только после создания комфортной связи
    переходят к сути. Без прочной социальной базы сделка не будет считаться
    надёжной.
\end{sloppypar}

\section{Задание 5. Понятие корпоративной культуры: совместная деятельность в политкультурной среде}
\begin{sloppypar}
Корпоративная культура — это система ценностей, норм, моделей поведения и общих
    правил, которые определяют, как сотрудники взаимодействуют друг с другом и
    с внешним миром внутри организации. Она задаёт атмосферу работы, стиль
    управления, принципы принятия решений, уровень формальности, отношение к
    инновациям, конфликтам и ответственности. Это не набор деклараций, а
    реальный «дух организации», который проявляется в ежедневной практике.

    Когда корпоративная культура формируется в политкультурной среде — то есть
    в коллективе, где работают представители разных стран, этнических групп и
    культурных традиций — её значение возрастает многократно. Такая среда
    требует согласования различных норм и стилей общения. Работники могут
    по-разному воспринимать авторитет, субординацию, индивидуальную и
    коллективную ответственность, сроки, дистанцию общения и способы выражения
    эмоций. Без единого корпоративного каркаса это быстро приводит к
    недопониманию.

    Корпоративная культура в политкультурной среде выполняет объединяющую
    функцию. Она становится общим кодексом поведения, который позволяет
    сотрудникам из разных культур работать эффективно и предсказуемо. Единые
    ценности — уважение, прозрачность, командная работа, ответственность —
    служат опорой, а стандартизированные правила уменьшают риски конфликтов.

    Кроме того, такая культура должна учитывать культурные различия и не
    навязывать одной группе доминирующую модель. Эффективная корпоративная
    культура в многонациональной среде опирается на принципы инклюзивности:
    прозрачные процедуры, уважение к разнообразию, гибкие подходы к
    коммуникации и признание различных способов решения задач.
\end{sloppypar}


\end{document}
